% ===== overview =====
\begin{frame}{이번 주차 개요}
  \begin{itemize}
    \item \kb{7.1 n-step TD} --- MC와 TD(0) 사이의 다이얼, U자 곡선
    \item \kb{7.2 적격흔적} --- TD($\lambda$)로 모든 상태를 동시에 갱신
    \item \kb{7.3 Dyna-Q} --- 모델 저장 + 가상 갱신으로 약 9.5배 가속
  \end{itemize}
  \vskip0.8em
  \begin{block}{이전·다음과의 연결}
    n-step TD, TD($\lambda$), Dyna-Q는 모두 \textbf{``부트스트래핑을 얼마나
    할 것인가''}라는 하나의 질문에 대한 서로 다른 답. Week 8 캡스톤에서
    Gymnasium의 표 기반 환경으로 세 개를 직접 구현·비교하고,
    n-step 리턴은 Week 9 DQN의 기본 옵션으로 다시 만난다.
  \end{block}
\end{frame}
